\subsection{Scalar norms with genuine nonetale local multiplicities}
\label{subsec:local-algebra-norms}

The local calculation retains every composition factor. Nilpotent layers
contribute local module length, and inseparable residue extensions contribute
their full field degree. The product formula below requires an actual product
isomorphism into local algebras. A construction of that isomorphism for every
Artinian algebra is a separate producer, and no arbitrary finite-fibre norm
classification is claimed in this section.

\begin{theorem}[Determinants on actual local-module composition factors]
\label{spectrum:local-module-determinant}
\lean{EnvelopeProperty.det_lsmul_simple_local,
EnvelopeProperty.det_lsmul_eq_residueNorm_pow_length}
\leanok
\uses{spectrum:algebra-norm-base-change}
Let $E$ be a commutative local algebra over any field $K$ and let $M$
be an $E$-module of finite length, finite-dimensional over $K$ with compatible
scalar actions. Multiplication by $b\in E$ on $M$ has determinant
\[
 \det_K(b\mid M)=N_{\kappa(E)/K}(\overline b)^{\ell_E(M)}.
\]
\end{theorem}
\begin{proof}
An actual simple $E$-module is isomorphic to a quotient by a maximal ideal.
Locality identifies that ideal with the unique maximal ideal, so the module
is genuinely isomorphic to $\kappa(E)$. Restrict the module isomorphism to
$K$ and conjugate scalar action: the operator of $b$ becomes multiplication
by $\overline b$ on this residue field. Its determinant is the residue-field
norm.

Induct on the genuine finite-length module construction. The zero module
has determinant one and length zero. In a simple quotient extension
$0\to N\to M\to Q\to0$, the $E$-submodule $N$ is invariant under $b$.
All underlying $K$-spaces are finite-dimensional, so the invariant-submodule
basis and quotient basis give a block triangular multiplication matrix.
Its determinant is the product of the determinants on $N$ and $Q$.
The induction hypothesis treats $N$ and the simple-module computation treats
$Q$. Actual module length is additive in the exact sequence, with
$\ell_E(Q)=1$, proving the exponent. This argument counts each composition
factor and uses no reducedness or separability.
\end{proof}

\begin{theorem}[A finite local algebra's scalar norm and degree]
\label{spectrum:local-algebra-norm-length}
\lean{EnvelopeProperty.algebraNorm_eq_residueNorm_pow_length,
EnvelopeProperty.finrank_eq_length_mul_residue_finrank}
\leanok
\uses{spectrum:local-module-determinant}
For every finite-dimensional commutative local $K$-algebra $E$,
\[
 N_{E/K}(b)=N_{\kappa(E)/K}(\overline b)^{\ell_E(E)},
 \qquad
 \dim_K E=\ell_E(E)[\kappa(E):K].
\]
\end{theorem}
\begin{proof}
Finiteness over the field makes $E$ Noetherian and Artinian as an
$E$-module, hence of finite length. Its multiplication determinant is its
algebra norm, so the preceding theorem applied to the regular module gives
the first identity.

For the second, the actual local-ring length restriction theorem gives
\[
 \ell_K(E)=\ell_E(E)\ell_{\kappa(K)}(\kappa(E)).
\] The canonical map
$K\to\kappa(K)$ is surjective because $K$ is already a field, so restriction
along this surjection preserves the residue-module length. Length of a
finite-dimensional vector space is its dimension. Taking finite natural
lengths yields the stated equality. Thus the local exponent times the
actual residue degree recovers the whole dimension, including nilpotent
layers and inseparable residue extensions.
\end{proof}

\begin{theorem}[Complete-field scalar norms including inseparability]
\label{spectrum:inseparable-scalar-field-norm}
\lean{EnvelopeProperty.valuation_algebraNorm_eq_pow_of_finite}
\leanok
\uses{spectrum:complete-norm-extension}
Let $H$ be a complete ultrametric normed field, $L/H$ any finite field
extension, and let $w$ explicitly extend the norm of $H$. Then
\[
 |N_{L/H}(b)|=w(b)^{[L:H]}.
\]
No separability or infinitude hypothesis is needed.
\end{theorem}
\begin{proof}
Extend the norm of $H$ to an algebraic closure $\Omega$. Let $F=H(b)$.
The genuine power basis of $F$ identifies each root of the minimal
polynomial of $b$ in $\Omega$ with an $H$-embedding $F\to\Omega$ taking its
generator to that root. Complete-field norm uniqueness identifies every
root value with $w(b)$.

The exact scalar norm formula is the product of the minimal-polynomial
root multiset raised to $[L:F]$. Preserve that multiset with all of its
repetitions: even for inseparable polynomials its cardinality, after
splitting in $\Omega$, is the minimal-polynomial degree $[F:H]$.
Multiplicativity thus gives $w(b)^{[F:H][L:F]}$, and genuine tower
dimension gives $[L:H]$. This computation explicitly retains the
inseparable multiplicities rather than counting distinct embeddings of $L$.
\end{proof}

\begin{theorem}[Valued local factors and their actual multiplicities]
\label{spectrum:nonetale-local-valued-norm}
\lean{EnvelopeProperty.valuation_algebraNorm_eq_pow_of_local,
EnvelopeProperty.valuation_algebraNorm_eq_prod_of_localAlgEquiv}
\leanok
\uses{spectrum:local-algebra-norm-length,spectrum:inseparable-scalar-field-norm}
Let $E$ be finite local over a complete ultrametric field $H$, and prescribe
an extending valuation $w$ on its actual residue field. Then
\[
 |N_{E/H}(b)|=w(\overline b)^{\ell_E(E)[\kappa(E):H]}.
\]
If a commutative $H$-algebra has an actual isomorphism
$E\simeq\prod_i E_i$ into finite local $H$-algebras and each actual
$\kappa(E_i)$ carries its extending valuation $w_i$, then
\[
 |N_{E/H}(b)|=\prod_i
 w_i(\overline{b_i})^{\ell_{E_i}(E_i)[\kappa(E_i):H]}.
\]
\end{theorem}
\begin{proof}
For a local factor, the algebraic norm-length identity replaces its
norm by the residue-field norm raised to its local length. Apply the
preceding complete-field norm identity to the residue field, retaining
its full degree, and multiply the two exponents. This proves the local
formula with both load-bearing multiplicities.

For the product statement, transport the actual multiplication operator
through the supplied algebra isomorphism. Norm invariance and the norm of
a finite product express its scalar norm as the product of the local
factor norms. Apply the local formula to every factor and use valuation
multiplicativity. The hypothesis is a genuine whole-algebra isomorphism,
not a scalar norm identity or a point-value correspondence. No reducedness,
etaleness, separability or infinitude is imposed, and existence of a
particular global decomposition is outside this statement.
\end{proof}
